\scanpage{89}
% archive foliation: 88
% lines 1--4
\uline{Groupes prolibres} spéciaux.
\[
\left\{\begin{aligned}
&\pi\text{ groupe profini}\\
&T\quad\widehat{\mathbb Z}\text{ module libre }\deleted{\simeq\widehat{\mathbb Z}\text{ (isomorphisme \unsure{non} précisé)}}\inserted{\text{de rang }1}\\
&\Phi_i\in\operatorname{Hom\,ext}(L,\pi)\deleted{\unsure{\pi}}\qquad 0\leq i\leq\nu\qquad(\nu\geq2)
\end{aligned}\right.
\]
% lines 5--9
\uuline{Hyp} :
\[
\left\{\begin{aligned}
&\exists e\in T^*\quad\text{base de }L\\
&(\varphi_i)_{0\leq i\leq\nu}\in\prod_{0\leq i\leq\nu}\Phi_i
\end{aligned}\right\}\quad
\begin{gathered}
\text{tel que, posant }h_i=\varphi_i(e),\\
\text{l'hom de}
\end{gathered}
\]
\diagram{\begin{tikzcd}[ampersand replacement=\&]
\inserted{\widehat L_\nu}\deleted{\unsure{\widehat F([0,\nu])}}\arrow[r]\&\pi\\
\lambda_i\arrow[r,mapsto]\&h_i
\end{tikzcd}}
\inserted{groupe prolibre \deleted{sur $[0,\nu]$}
\unsure{engendré par} $\lambda_0,\ldots,\lambda_\nu$ avec $\lambda_0\cdots\lambda_\nu=1$}
soit un \uline{isomorphisme}
% lines 10--14
En d'autres termes
\[
\begin{aligned}
\text{Soit }\pi_0&\deleted{\unsure{\Leftarrow|\leftarrow}\text{\unsure{groupe pro.}}}\inserted{=\widehat L_\nu=}\ (\lambda_0,\ldots,\lambda_\nu\mid\lambda_0\lambda_1\cdots\lambda_\nu=1)^{\wedge}\\
T_0&=\widehat{\mathbb Z},\ e_0=1
\end{aligned}
\]
$\Phi_{0i}$ classe de
\diagram{\begin{tikzcd}[ampersand replacement=\&]
\widehat{\mathbb Z}\arrow[r]\&\unsure{\pi_\nu|\pi_0}\\
e_0\arrow[r,mapsto]\&\lambda_i
\end{tikzcd}}
% lines 15--17
on veut que $(\pi,T,(\Phi_i))$ soit isom. à $(\pi_0,T_0,(\Phi_{0i}))$
\[
A_\nu=\operatorname{Aut}(\widehat L_\nu,T_0,(\Phi_{0i}))\xrightarrow{\chi_\nu}
\underset{\substack{\Vert\\\operatorname{Aut}(T_0)}}{\widehat{\mathbb Z}^{*}}
\]
\margin{on prouve que c'est surjectif}
% lines 18--20
\diagram{\begin{tikzcd}[ampersand replacement=\&]
\widehat L_\nu\arrow[r,"i"]\&\unsure{A.\operatorname{Ker}\chi_\nu|A=\operatorname{Ker}\chi_\nu}\subset A_\nu\\
g\arrow[r,mapsto]\&(\operatorname{int}(g),\operatorname{id}_{T_0})
\end{tikzcd}}
injectif
\[
u\,i(g)\,u^{-1}=i(u(g))
\]
i.e $\widehat L_\nu$ invariant ds $A_\nu$
% line 21
$A_\nu/\widehat L_\nu=\Gamma_\nu$\quad on veut élucider sa structure.
